% =============================================================================
% Lennard-Jones Generalization Test: v2 Results
% =============================================================================
% Standalone results document for the LJ PES experiments.
% Informed by LITERATURE_REVIEW_AND_SYNTHESIS.tex but kept separate.
% =============================================================================

\documentclass[11pt,a4paper]{article}

\usepackage[utf8]{inputenc}
\usepackage[T1]{fontenc}
\usepackage{amsmath,amssymb,amsthm}
\usepackage{mathtools}
\usepackage{bm}
\usepackage{booktabs}
\usepackage{longtable}
\usepackage{hyperref}
\usepackage[margin=1in]{geometry}
\usepackage{enumitem}
\usepackage{xcolor}
\usepackage{multirow}
\usepackage{colortbl}
\usepackage{tcolorbox}
\usepackage{tabularx}

% --- Color palette ---
\definecolor{ljblue}{HTML}{1565C0}
\definecolor{ljlight}{HTML}{E3F2FD}
\definecolor{dftborange}{HTML}{E65100}
\definecolor{dftblight}{HTML}{FFF3E0}
\definecolor{successgreen}{HTML}{2E7D32}
\definecolor{successlight}{HTML}{E8F5E9}
\definecolor{warnred}{HTML}{C62828}
\definecolor{warnlight}{HTML}{FFEBEE}
\definecolor{keybox}{HTML}{F3E5F5}
\definecolor{keyborder}{HTML}{7B1FA2}

% --- Highlighted-box environments ---
\tcbset{
  keyresult/.style={
    colback=keybox, colframe=keyborder,
    fonttitle=\bfseries\large, coltitle=white, colbacktitle=keyborder,
    boxrule=1.5pt, arc=4pt, left=8pt, right=8pt, top=6pt, bottom=6pt,
  },
  finding/.style={
    colback=successlight, colframe=successgreen,
    fonttitle=\bfseries, coltitle=white, colbacktitle=successgreen,
    boxrule=1pt, arc=3pt, left=6pt, right=6pt, top=4pt, bottom=4pt,
  },
  warning/.style={
    colback=warnlight, colframe=warnred,
    fonttitle=\bfseries, coltitle=white, colbacktitle=warnred,
    boxrule=1pt, arc=3pt, left=6pt, right=6pt, top=4pt, bottom=4pt,
  },
}

\newcommand{\vect}[1]{\bm{#1}}
\newcommand{\mat}[1]{\mathbf{#1}}
\newcommand{\norm}[1]{\lVert #1 \rVert}
\newcommand{\abs}[1]{\lvert #1 \rvert}
\newcommand{\hlnum}[1]{{\Large\bfseries\color{ljblue}#1}}
\newcommand{\hldftb}[1]{{\large\bfseries\color{dftborange}#1}}

\theoremstyle{definition}
\newtheorem{remark}{Remark}

\title{Newton--Raphson Minimization on Lennard-Jones:\\
  Generalization of the v13 Optimizer to a Pairwise Potential}
\author{Transition-State Tools Project}
\date{March 2026}

\begin{document}
\maketitle

% =============================================================================
\section{Introduction}
% =============================================================================

The v13 Newton--Raphson optimizer achieves 93--98\% strict convergence on
DFTB0 across noise levels 0.5--2.0\,\AA{} (Transition1x, $N{=}287$
samples).  Does this generalize beyond DFTB0?

We swap DFTB0 for a \textbf{pairwise Lennard-Jones} (LJ) potential---same
geometries, same noise, same optimizer, same criteria---and find:

\begin{tcolorbox}[keyresult, title=Central Result]
\begin{itemize}[nosep,leftmargin=*]
  \item The optimizer \textbf{works}: \hlnum{95--99\%} relaxed convergence
    at every noise level
  \item Strict convergence \textbf{collapses}: \hlnum{92\% $\to$ 7\%} as
    noise increases---entirely due to \textbf{ghost modes} from LJ's
    missing angular terms, not optimizer failure
  \item Under relaxed criterion, LJ (\hlnum{97.5\%} avg) actually
    \textbf{exceeds} DFTB0 (\hldftb{96.1\%} avg)
\end{itemize}
\end{tcolorbox}

\paragraph{Why LJ?}  Three properties make it an ideal stress test:
(1)~exact analytical derivatives (no SCF noise),
(2)~fundamentally different PES topology (purely pairwise, no angular terms),
(3)~${\sim}1000\times$ faster, enabling rapid grid sweeps.
LJ uses UFF parameters with $\sigma/3$ scaling to match covalent bond
lengths (Appendix~\ref{app:lj-potential}).

% =============================================================================
\section{Setup}
\label{sec:setup}
% =============================================================================

The full optimizer (``v12b'': RFO step builder, polynomial line search,
Eckart projection, three kick mechanisms) and noise model are detailed in
Appendices~\ref{app:optimizer}--\ref{app:noise}.  Here we summarize the
key parameters and grid design.

\subsection{Grid design}

15 combinations: 3~configs $\times$ 4~noise levels $+$ 3~ablations at
$n{=}2.0$\,\AA.

\begin{table}[h]
\centering\small
\begin{tabular}{clccl}
\toprule
\# & Configuration & Steps & Noise & Purpose \\
\midrule
1 & RFO+PLS baseline       & 20k & all       & Same-seed control \\
2 & RFO+PLS baseline       & 50k & all       & Step-count reference \\
3 & Full v12b (all kicks)  & 50k & all       & Main experiment \\
\midrule
4 & Full v12b (all kicks)  & 20k & 2.0 only  & Kick effect control \\
5 & Adapt+probe (no late)  & 50k & 2.0 only  & Does 50k replace late? \\
6 & Probe+late (no adapt)  & 50k & 2.0 only  & Is adaptive needed? \\
\bottomrule
\end{tabular}
\caption{LJ v2 grid.  All: $N{=}287$ Transition1x samples, seed\,=\,42,
  48 workers $\times$ 4 threads.}
\label{tab:grid}
\end{table}

\subsection{Convergence criteria}

\begin{description}[nosep]
  \item[\color{successgreen}Strict:] $\norm{\vect{f}}_\text{mean}
    < \varepsilon$ \textbf{and} $n_\text{neg} = 0$ \quad
    (no negative vibrational eigenvalues)
  \item[\color{ljblue}Relaxed:] $\norm{\vect{f}}_\text{mean}
    < \varepsilon$ \textbf{and} $\lambda_\text{min} \geq -0.01$ \quad
    (eigenvalues within $\sqrt{\varepsilon}$ of zero are accepted)
\end{description}

\noindent
The relaxed threshold $\sqrt{\varepsilon} = \sqrt{10^{-4}} = 0.01$ follows
from the $\varepsilon^{1/2}$ second-order optimality theorem
(Boumal et al.): an RFO-type method guarantees
$\lambda_\text{min} \geq -C\sqrt{\varepsilon}$.

\subsection{Key parameters}
\label{sec:params}

\begin{center}\small
\begin{tabular}{lll}
\toprule
Parameter & Value & Purpose \\
\midrule
Calculator & LJ (UFF, $\sigma/3$) & PES evaluation \\
Force convergence $\varepsilon$ & $10^{-4}$\,eV/\AA & first-order criterion \\
Relaxed threshold $\tau$ & 0.01 & second-order criterion \\
Max atom displacement & 1.3\,\AA & initial trust radius \\
Trust radius floor & 0.01\,\AA & prevents TR collapse \\
Min interatomic distance & 0.5\,\AA & rejects unphysical steps \\
\bottomrule
\end{tabular}
\end{center}

\noindent All parameters identical to DFTB0 v13 except the calculator.
Full parameter table in Appendix~\ref{app:params}.

% =============================================================================
% =============================================================================
\section{Results}
\label{sec:results}
% =============================================================================
% =============================================================================

% -------------------------------------------------------------------
\subsection{The headline: strict convergence collapses on LJ}
% -------------------------------------------------------------------

\begin{table}[h]
\centering
\caption{Strict convergence ($n_\text{neg}{=}0$) out of $N{=}287$.
  \colorbox{warnlight}{\strut Red shading} = below 50\%.}
\label{tab:lj-strict}
\begin{tabular}{l
  >{\columncolor{white}}c
  >{\columncolor{white}}c
  >{\columncolor{white}}c
  >{\columncolor{white}}c}
\toprule
Configuration & 0.5\,\AA & 1.0\,\AA & 1.5\,\AA & 2.0\,\AA \\
\midrule
RFO+PLS 20k & 257 (89.5\%) & 168 (58.5\%)
  & \cellcolor{warnlight} 70 (24.4\%)
  & \cellcolor{warnlight} 13 (4.5\%) \\
RFO+PLS 50k & 258 (89.9\%) & 180 (62.7\%)
  & \cellcolor{warnlight} 92 (32.1\%)
  & \cellcolor{warnlight} 17 (5.9\%) \\
Full v12b 50k & \textbf{265} (92.3\%) & \textbf{185} (64.5\%)
  & \cellcolor{warnlight} \textbf{96} (33.4\%)
  & \cellcolor{warnlight} \textbf{19} (6.6\%) \\
\midrule
\rowcolor{dftblight}
\textit{DFTB0 v13 (ref)} & \textit{281 (97.9\%)} & \textit{278 (96.9\%)}
  & \textit{277 (96.5\%)} & \textit{267 (93.0\%)} \\
\bottomrule
\end{tabular}
\end{table}

Strict convergence on LJ drops from 92\% to \textbf{7\%} as noise
increases.  On DFTB0, it drops only from 98\% to 93\%.
\emph{But this is not an optimizer failure}---as the next table shows.

% -------------------------------------------------------------------
\subsection{The real story: relaxed convergence is excellent}
% -------------------------------------------------------------------

\begin{table}[h]
\centering
\caption{Relaxed convergence ($\lambda_\text{min} \geq -0.01$).
  \colorbox{successlight}{\strut Green shading} = above 97\%.}
\label{tab:lj-relaxed}
\begin{tabular}{l
  >{\columncolor{white}}c
  >{\columncolor{successlight}}c
  >{\columncolor{successlight}}c
  >{\columncolor{successlight}}c}
\toprule
Configuration & 0.5\,\AA & 1.0\,\AA & 1.5\,\AA & 2.0\,\AA \\
\midrule
RFO+PLS 20k & 273 (95.1\%) & 279 (97.2\%) & 281 (97.9\%) & 284 (99.0\%) \\
RFO+PLS 50k & 273 (95.1\%) & 279 (97.2\%) & 282 (98.3\%) & 285 (99.3\%) \\
Full v12b 50k & 273 (95.1\%) & 279 (97.2\%) & 282 (98.3\%) & \textbf{285 (99.3\%)} \\
\bottomrule
\end{tabular}
\end{table}

\begin{tcolorbox}[finding, title=Finding 1: The optimizer works]
Under relaxed convergence, LJ achieves \hlnum{95--99\%} at every noise
level.  At $n{=}2.0$, LJ relaxed (\hlnum{99.3\%}) \textbf{exceeds}
DFTB0 strict (\hldftb{93.0\%}).  The optimizer finds force-converged
minima on a fundamentally different PES---the ``failures'' are a
\emph{criterion problem}, not an optimizer problem.
\end{tcolorbox}

% -------------------------------------------------------------------
\subsection{The strict--relaxed gap: the central diagnostic}
% -------------------------------------------------------------------

\begin{table}[h]
\centering
\caption{The gap between strict and relaxed convergence.  LJ's gap
  explodes with noise; DFTB0's is negligible.}
\label{tab:gap}
\begin{tabular}{c
  >{\color{ljblue}\bfseries}c
  c
  >{\color{ljblue}\bfseries}c
  c}
\toprule
Noise & LJ strict & LJ relaxed & \textnormal{LJ gap} & DFTB0 gap \\
\midrule
0.5\,\AA & 92.3\% & 95.1\% & 2.8\,pts & ${\sim}$0\,pts \\
1.0\,\AA & 64.5\% & 97.2\% & 32.7\,pts & ${\sim}$0\,pts \\
1.5\,\AA & 33.4\% & 98.3\% & \cellcolor{warnlight}64.9\,pts & ${\sim}$0\,pts \\
2.0\,\AA & 6.6\%  & 99.3\% & \cellcolor{warnlight}92.7\,pts & ${\sim}$0\,pts \\
\bottomrule
\end{tabular}
\end{table}

The gap grows from 3\,pts to \textbf{93\,pts} as noise increases.  On
DFTB0 it is always ${\sim}0$.  Why?

% -------------------------------------------------------------------
\subsection{Aggregate: LJ vs.\ DFTB0 v13}
\label{sec:aggregate}
% -------------------------------------------------------------------

\begin{table}[h]
\centering
\caption{Aggregate comparison (full v12b 50k, mean over 4 noise levels).}
\label{tab:aggregate}
\begin{tabular}{l
  >{\columncolor{ljlight}}c
  >{\columncolor{dftblight}}c
  >{\columncolor{ljlight}}c
  >{\columncolor{dftblight}}c}
\toprule
 & \multicolumn{2}{c}{\textbf{Strict} ($n_\text{neg}{=}0$)}
 & \multicolumn{2}{c}{\textbf{Relaxed} ($\lambda_\text{min}{\geq}{-}0.01$)} \\
\cmidrule(lr){2-3} \cmidrule(lr){4-5}
Metric
  & \cellcolor{ljlight}\textbf{\color{ljblue}LJ}
  & \cellcolor{dftblight}\textbf{\color{dftborange}DFTB0}
  & \cellcolor{ljlight}\textbf{\color{ljblue}LJ}
  & \cellcolor{dftblight}\textbf{\color{dftborange}DFTB0} \\
\midrule
Mean rate & 49.2\% & 96.1\% & \textbf{97.5\%} & 96.6\% \\
Total converged & 565/1119 & 1103/1148 & \textbf{1119/1119} & ${\sim}$1107/1148 \\
Spread & 85.7\,pts & 4.9\,pts & \textbf{4.2\,pts} & ${\sim}$4.9\,pts \\
\bottomrule
\end{tabular}
\end{table}

\begin{tcolorbox}[finding, title=Finding 2: LJ exceeds DFTB0 under relaxed criterion]
Under strict criterion, LJ averages 49\% vs.\ DFTB0's 96\%.
Under relaxed criterion, LJ (\textbf{97.5\%}) \emph{exceeds} DFTB0
(96.6\%), and the noise-level spread is \emph{smaller} (4.2 vs.\ 4.9\,pts).
The optimizer is equally robust on both PES once ghost modes are excluded.
\end{tcolorbox}

% =============================================================================
\section{Why Strict Convergence Collapses: Ghost Modes}
\label{sec:ghost}
% =============================================================================

\subsection{What are ghost modes?}

A \textbf{ghost mode} is a vibrational eigenvalue that is negative but
negligible ($|\lambda| < 10^{-3}$).  The PES curvature in that direction
is \emph{effectively zero}; the sign is determined by floating-point
rounding after Eckart projection.

\begin{tcolorbox}[warning, title=Root cause: LJ has no angular terms]
\textbf{DFTB0} has angular contributions via its tight-binding Hamiltonian
$\Rightarrow$ positive curvature in all $3N{-}6$ directions
$\Rightarrow$ ghost modes $= 1\%$ of failures.

\medskip

\textbf{LJ} is purely pairwise: $V = \sum_{i<j} V_\text{LJ}(r_{ij})$.
Directions that change bond angles without changing any distance have
\emph{exactly zero} Hessian curvature $\Rightarrow$ ghost modes $= 67\%$
of failures at $n{=}2.0$.
\end{tcolorbox}

\subsection{The coin-flip collapse}

If the Hessian has $k$ ghost eigenvalues (signs effectively random), the
probability of passing strict ($n_\text{neg}{=}0$) is:
\[
  P(\text{strict}) \approx \left(\tfrac{1}{2}\right)^k
\]

As noise increases, more directions become flat ($k$ grows), and strict
convergence collapses \emph{exponentially}:

\begin{center}
\begin{tabular}{cccl}
\toprule
Noise & Approx.\ $k$ & $(1/2)^k$ & Observed strict \\
\midrule
0.5\,\AA & 0--1 & 50--100\% & 92.3\% \\
1.0\,\AA & 1--2 & 25--50\% & 64.5\% \\
1.5\,\AA & 2--3 & 12--25\% & 33.4\% \\
2.0\,\AA & 3--5 & 3--12\%  & 6.6\%  \\
\bottomrule
\end{tabular}
\end{center}

The observed rates match the coin-flip model quantitatively.  DFTB0 has
$k \approx 0$ at all noise levels, so $(1/2)^0 = 1$ and strict $\approx$
relaxed.

\begin{tcolorbox}[finding, title=Finding 3: The collapse is a criterion problem]
Strict convergence ($n_\text{neg}{=}0$) demands positive curvature in
directions where the PES is \emph{genuinely flat}.  This is not an
optimizer failure---it is an \textbf{inappropriate criterion} for pairwise
potentials.  The relaxed criterion resolves this completely.
\end{tcolorbox}

\subsection{Threshold quality: is 0.01 the right value?}

The cascade data (Appendix~\ref{app:cascade}) shows that convergence
\textbf{plateaus by threshold 0.001--0.002}---the 0.01 bound from the
$\varepsilon^{1/2}$ theorem is $5\times$ looser than necessary.

\begin{center}\small
\begin{tabular}{ccccc}
\toprule
Noise & Plateau thr & Plateau value & Rate excl.\ errors \\
\midrule
0.5\,\AA & 0.002 & 273/273 & \textbf{100\%} \\
1.0\,\AA & 0.002 & 279/279 & \textbf{100\%} \\
1.5\,\AA & 0.001 & 279/282 & 98.9\% \\
2.0\,\AA & 0.002 & 284/285 & 99.6\% \\
\bottomrule
\end{tabular}
\end{center}

There is a \textbf{clean gap}: no sample has an eigenvalue between $-0.01$
and $-0.002$.  Ghost modes cluster in $(-10^{-3}, 0)$; genuine saddle
eigenvalues would be much more negative.  The relaxed threshold is safe.

% =============================================================================
\section{LJ vs.\ DFTB0: Key Contrasts}
\label{sec:comparison}
% =============================================================================

\renewcommand{\arraystretch}{1.3}
\begin{table}[h]
\centering
\caption{Side-by-side comparison (full v12b 50k).}
\label{tab:contrasts}
\begin{tabular}{p{3.8cm}
  >{\columncolor{ljlight}\centering\arraybackslash}p{4.5cm}
  >{\columncolor{dftblight}\centering\arraybackslash}p{4.5cm}}
\toprule
& \textbf{\color{ljblue}LJ} & \textbf{\color{dftborange}DFTB0 v13} \\
\midrule
Strict range & 92\% $\to$ 7\% & 98\% $\to$ 93\% \\
Relaxed range & 95\% $\to$ 99\% & $\sim$98\% $\to$ $\sim$93\% \\
Dominant failure & Ghost modes (67\%) & Oscillation (8.2\%) \\
Kick gain at $n{=}2.0$ & $+$0.7\,pts & $+$12.2\,pts \\
Median steps ($n{=}2.0$) & 817 & $\sim$11\,000 \\
Mean wall ($n{=}2.0$) & 11\,s & 368\,s \\
Ablation spread & 5 samples (noise) & 35 samples (meaningful) \\
Hardest samples & \{106, 35, 178\} & \{136, 150, 209\} \\
\bottomrule
\end{tabular}
\end{table}
\renewcommand{\arraystretch}{1.0}

\begin{tcolorbox}[finding, title=Finding 4: Performance is PES-specific]
\begin{enumerate}[nosep]
  \item \textbf{Kicks help less on LJ} ($+$0.7\,pts vs.\ $+$12.2\,pts at
    $n{=}2.0$) because ghost modes cannot be resolved by perturbation
  \item \textbf{LJ is 3--14$\times$ faster in steps, 9--33$\times$ in
    wall time}---the smooth PES is easier to optimize
  \item \textbf{Zero overlap in hardest samples}---difficulty is
    PES-specific, not geometry-specific
  \item \textbf{Ablation differences are noise on LJ} (5 samples) vs.\
    meaningful on DFTB0 (35 samples)
\end{enumerate}
\end{tcolorbox}

% =============================================================================
\section{Conclusions}
\label{sec:conclusions}
% =============================================================================

\begin{tcolorbox}[keyresult, title=Summary of Findings]
\begin{enumerate}[leftmargin=*]
  \item \textbf{The v13 optimizer generalizes to LJ.}  Relaxed
    convergence of \hlnum{95--99\%} across all noise levels on a
    fundamentally different PES.

  \item \textbf{Strict convergence is inappropriate for pairwise
    potentials.}  The collapse from 92\% to 7\% is entirely due to ghost
    modes from missing angular terms---not optimizer failure.

  \item \textbf{The $\bm{\varepsilon^{1/2}}$ theorem is validated.}
    The relaxed criterion $\lambda_\text{min} \geq -\sqrt{\varepsilon}$
    captures essentially all converged samples.

  \item \textbf{Ghost modes dominate LJ (67\%); oscillation dominates
    DFTB0 (8.2\%).}  Remaining challenges are PES-specific.

  \item \textbf{LJ is 9--33$\times$ faster in wall time.}  The smooth
    PES confirms DFTB0's difficulty is intrinsic to its topology.

  \item \textbf{Kicks help less on LJ.}  Ghost modes cannot be resolved
    by perturbation; only a criterion change helps.
\end{enumerate}
\end{tcolorbox}

\subsection{Implications for the optimizer}

The LJ results argue strongly for \textbf{adopting the relaxed convergence
criterion as the primary metric} on any PES.  Strict convergence is
suitable only when the PES is known to have well-defined curvature in all
vibrational directions, and even then the $\varepsilon^{1/2}$ theorem
says eigenvalues in $[-\sqrt{\varepsilon}, 0)$ should be accepted.

For future work on potentials with flat directions (force fields, classical
pair potentials, coarse-grained models), the relaxed criterion should be
used from the start.

% #############################################################################
%
%                            A P P E N D I C E S
%
% #############################################################################
\clearpage
\appendix

% =============================================================================
\section{Lennard-Jones Potential}
\label{app:lj-potential}
% =============================================================================

The LJ potential is a pairwise sum:
\begin{equation}
  V = \sum_{i<j} 4\,\varepsilon_{ij}
      \left[\left(\frac{\sigma_{ij}}{r_{ij}}\right)^{12}
            - \left(\frac{\sigma_{ij}}{r_{ij}}\right)^{6}\right],
  \label{eq:lj}
\end{equation}
with analytical gradient and Hessian:
\begin{align}
  \frac{\partial V}{\partial r_{ij}}
    &= \frac{4\varepsilon_{ij}}{r_{ij}}
       \left[-12\left(\frac{\sigma_{ij}}{r_{ij}}\right)^{12}
             + 6\left(\frac{\sigma_{ij}}{r_{ij}}\right)^{6}\right],
  \label{eq:lj-deriv} \\[4pt]
  \frac{\partial^2 V}{\partial r_{ij}^2}
    &= \frac{4\varepsilon_{ij}}{r_{ij}^2}
       \left[156\left(\frac{\sigma_{ij}}{r_{ij}}\right)^{12}
             - 42\left(\frac{\sigma_{ij}}{r_{ij}}\right)^{6}\right].
  \label{eq:lj-hess-rr}
\end{align}
The Cartesian Hessian off-diagonal $3{\times}3$ blocks are:
\begin{equation}
  H_{ia,jb} = -\!\left(
    \frac{\partial^2 V}{\partial r^2} - \frac{1}{r}\frac{\partial V}{\partial r}
  \right) \hat{r}_a \hat{r}_b
  \;-\; \frac{1}{r}\frac{\partial V}{\partial r}\;\delta_{ab},
  \quad i \ne j,
  \label{eq:lj-hess-cart}
\end{equation}
with diagonal blocks $H_{ii} = -\sum_{j \ne i} H_{ij}$.

\paragraph{UFF parameters and $\sigma$-scaling.}
Per-atom $\sigma_i, \varepsilon_i$ from UFF
(Rapp\'{e} et al., \textit{J.\ Am.\ Chem.\ Soc.}\ 1992), with
Lorentz--Berthelot mixing:
$\sigma_{ij} = (\sigma_i + \sigma_j)/2$,
$\varepsilon_{ij} = \sqrt{\varepsilon_i \varepsilon_j}$.
Raw UFF $\sigma \in [2.6, 3.7]$\,\AA{} ($r_\text{eq} = 2^{1/6}\sigma
\approx 2.9$--$4.2$\,\AA).  We scale by $1/3$
(\texttt{lj\_sigma\_scale\,=\,0.3333}) to shrink equilibria to
${\sim}1.0$--$1.3$\,\AA, matching covalent bond lengths.
$1\,\text{kcal/mol} = 0.0434\,\text{eV}$.

\paragraph{What LJ lacks.}
$V$ depends only on distances $r_{ij}$, not angles or dihedrals.
Directions changing angles without changing any $r_{ij}$ have
\emph{exactly zero} Hessian curvature---the origin of ghost modes
(\S\ref{sec:ghost}).

% =============================================================================
\section{The v12b Optimizer}
\label{app:optimizer}
% =============================================================================

The optimizer is a Newton--Raphson minimizer in Cartesian coordinates,
assembled incrementally across v1--v12b.

\subsection{RFO step builder (v10b)}

The raw Hessian is mass-weighted and projected to the $(3N{-}k)$-dimensional
vibrational subspace (Section~\ref{app:eckart}).  The reduced Hessian is
diagonalized:
\begin{equation}
  \mat{H}_\text{red} = \mat{Q}_\text{vib}^\top \mat{M}^{-1/2}\mat{H}\mat{M}^{-1/2} \mat{Q}_\text{vib},
  \quad
  \mat{H}_\text{red}\,\vect{v}_i = \lambda_i\,\vect{v}_i.
\end{equation}
The step uses the RFO augmented Hessian (Schlegel 2011):
\begin{equation}
  \begin{pmatrix} \mat{H}_\text{red} & \vect{g} \\
                  \vect{g}^\top & 0 \end{pmatrix}
  \begin{pmatrix} \vect{s} \\ 1 \end{pmatrix}
  = \mu \begin{pmatrix} \vect{s} \\ 1 \end{pmatrix},
\end{equation}
with $\mu$ found by Newton's method on the secular equation
$\sum_i c_i^2/(\lambda_i - \mu) + \mu = 0$.
Step coefficients: $\Delta x_i = -c_i / (\lambda_i - \mu)$.
For minimization, $\mu < \lambda_\text{min}$ $\Rightarrow$ guaranteed
descent.  RFO automatically interpolates between steepest descent and
Newton---no hyperparameters.

\subsection{Trust region and displacement capping}

The step is capped per-atom:
$\Delta\vect{x} \leftarrow \Delta\vect{x} \cdot R / \max_i\norm{\Delta\vect{x}_i}$
when $\max_i\norm{\Delta\vect{x}_i} > R$.  Trust radius updated via
$\rho = \Delta E_\text{actual} / \Delta E_\text{predicted}$:
\begin{equation}
  R \leftarrow \begin{cases}
    \min(1.5\,R,\;\text{MAD}) & \rho > 0.75 \text{ and at boundary} \\
    \max(0.25\,\norm{\Delta\vect{x}}_\text{max},\;R_\text{floor}) & \rho < 0.25 \\
    R & \text{otherwise}
  \end{cases}
\end{equation}
$R_\text{floor} = 0.01$\,\AA{} prevents collapse.  Steps violating
$d_\text{min} = 0.5$\,\AA{} are rejected.

\subsection{Polynomial line search (v10c)}

After each accepted step, a cubic $p(t)$ is fit on $[0,1]$ using energy and
directional derivatives at old/new positions.  If an interior minimum
$t^* \in (0.1, 0.9)$ with $p''(t^*) > 0$ exists and yields lower energy,
the step is refined to $t^*\,\Delta\vect{x}$.

\subsection{Eckart projection}
\label{app:eckart}

TR modes are removed \emph{by construction}, not threshold filtering.
An Eckart basis $\mat{B} \in \mathbb{R}^{3N \times k}$ ($k=5$ or $6$) is
built from mass-weighted translations and rotations.  The complementary
$\mat{Q}_\text{vib}$ gives the full-rank reduced Hessian
$\mat{H}_\text{red} = \mat{Q}_\text{vib}^\top \mat{H}_\text{mw} \mat{Q}_\text{vib}
\in \mathbb{R}^{(3N-k) \times (3N-k)}$---every eigenvalue is a genuine
vibrational frequency.  The gradient is also projected
(\texttt{--project-gradient-and-v}).

\subsection{Kick mechanisms (v12/v12b)}

Three perturbation mechanisms escape saddle points when $n_\text{neg} > 0$.
A \textbf{mode tracker} maintains eigenvector continuity across steps.

\paragraph{1.\ Oscillation kick (v12).}
Fires when energy oscillation is detected at the trust-radius floor for
$\geq 3$ consecutive checks.  Perturbs along the longest-stuck negative
eigenvector:
$\vect{x} \leftarrow \vect{x} + s \cdot \delta \cdot \vect{v}_\text{neg}$,
$s = \text{sign}(\vect{g} \cdot \vect{v}_\text{neg})$.
50-step cooldown.

\paragraph{2.\ Blind-mode kick with probe (v12b).}
Fires when a negative mode persists $\geq 100$ steps with gradient overlap
$< 0.1$ and force norm $< 0.1$\,eV/\AA.  Probes both $\pm\vect{v}$ at
$1, 2, 5, 10, 20 \times R_\text{floor}$; accepts first point where
$n_\text{neg}$ decreases.

\paragraph{3.\ Late escape (v12b).}
After step 15k, if $n_\text{neg} > 0$: displacement $\alpha = 0.1$\,\AA{}
along longest-stuck mode in both directions.  Accepted only if
$n_\text{neg}$ decreases.  500-step cooldown.

\paragraph{Adaptive kick scaling (v12b).}
$\delta = \min(C\sqrt{|\lambda_\text{neg}|},\; 10\,R_\text{floor})$,
$\delta \geq R_\text{floor}$, with $C = 0.1$.

% =============================================================================
\section{Starting Geometries and Noise Model}
\label{app:noise}
% =============================================================================

Starting geometry = midpoint:
$\vect{x}_0 = \tfrac{1}{2}(\vect{x}_R + \vect{x}_\text{TS})$.
Gaussian noise: $\vect{x}_\text{start} = \vect{x}_0 + \bm{\eta}$,
$\eta_{ia} \sim \mathcal{N}(0, \sigma^2)$, with
$\sigma \in \{0.5, 1.0, 1.5, 2.0\}$\,\AA.
Seed per sample: $\texttt{noise\_seed} + \texttt{sample\_index}$
(base seed = 42).

% =============================================================================
\section{Full Parameter Table}
\label{app:params}
% =============================================================================

\begin{center}\small
\begin{tabular}{lll}
\toprule
Parameter & Value & Purpose \\
\midrule
Calculator & LJ (UFF, $\sigma/3$) & PES evaluation \\
Force convergence $\varepsilon$ & $10^{-4}$\,eV/\AA & first-order criterion \\
Relaxed threshold $\tau$ & 0.01 & second-order criterion \\
Max atom displacement (MAD) & 1.3\,\AA & initial trust radius \\
Trust radius floor $R_\text{floor}$ & 0.01\,\AA & prevents TR collapse \\
Min interatomic distance & 0.5\,\AA & rejects unphysical steps \\
Noise seed & 42 & reproducibility \\
Late escape after & 15\,000 steps & last-resort saddle escape \\
Late escape $\alpha$ & 0.1\,\AA & late-escape displacement \\
Late escape cooldown & 500 steps & limits late-escape frequency \\
Adaptive kick $C$ & 0.1 & eigenvalue-scaled kick magnitude \\
Osc-kick patience & 3 & consecutive oscillation detections \\
Osc-kick cooldown & 50 steps & limits osc-kick frequency \\
Blind-kick patience & 100 steps & mode persistence before kick \\
Blind-kick overlap threshold & 0.1 & defines ``blind'' (low overlap) \\
Blind-kick force threshold & 0.1\,eV/\AA & kick only when forces are small \\
Blind-kick probe factors & $1,2,5,10,20$ & multiples of $R_\text{floor}$ \\
\bottomrule
\end{tabular}
\end{center}

% =============================================================================
\section{Cascade Tables}
\label{app:cascade}
% =============================================================================

\begin{table}[h]
\centering\small
\caption{Cascade: converged at each eigenvalue threshold (full v12b 50k
  and baselines).}
\label{tab:cascade}
\begin{tabular}{lccccccc}
\toprule
Config & 0.0 & 0.0001 & 0.0005 & 0.001 & 0.002 & 0.005 & 0.01 \\
\midrule
\multicolumn{8}{l}{\textit{n = 0.5\,\AA}} \\
\quad 20k baseline    & 257 & 264 & 266 & 269 & 271 & 273 & 273 \\
\quad 50k baseline    & 258 & 265 & 267 & 270 & 272 & 273 & 273 \\
\quad full v12b 50k   & 265 & 271 & 271 & 272 & 273 & 273 & 273 \\
\midrule
\multicolumn{8}{l}{\textit{n = 1.0\,\AA}} \\
\quad 20k baseline    & 168 & 252 & 261 & 264 & 267 & 279 & 279 \\
\quad 50k baseline    & 180 & 265 & 271 & 274 & 276 & 278 & 279 \\
\quad full v12b 50k   & 185 & 272 & 274 & 278 & 279 & 279 & 279 \\
\midrule
\multicolumn{8}{l}{\textit{n = 1.5\,\AA}} \\
\quad 20k baseline    & 70  & 218 & 249 & 259 & 268 & 280 & 281 \\
\quad 50k baseline    & 92  & 254 & 276 & 278 & 281 & 282 & 282 \\
\quad full v12b 50k   & 96  & 265 & 279 & 282 & 282 & 282 & 282 \\
\midrule
\multicolumn{8}{l}{\textit{n = 2.0\,\AA}} \\
\quad 20k baseline    & 13  & 196 & 260 & 272 & 277 & 282 & 284 \\
\quad 50k baseline    & 17  & 209 & 269 & 277 & 281 & 283 & 285 \\
\quad full v12b 50k   & 19  & 210 & 275 & 282 & 285 & 285 & 285 \\
\bottomrule
\end{tabular}
\end{table}

The knee is always at $0 \to 0.0001$: this is where the majority of ghost
modes live ($+$191 samples at $n{=}2.0$).

% =============================================================================
\section{Step Distributions and Computational Cost}
\label{app:steps}
% =============================================================================

\begin{table}[h]
\centering\small
\caption{Step and wall-time distributions (non-error samples).}
\begin{tabular}{lccccccc}
\toprule
Config & Min & Max & Mean & Median & P90 & Wall mean & Wall med \\
\midrule
\multicolumn{8}{l}{\textit{n = 0.5\,\AA}} \\
\quad 20k baseline    & 2 & 15\,080 & 1\,276 & 516 & 3\,224 & 5.3\,s & 1.4\,s \\
\quad 50k baseline    & 2 & 31\,986 & 1\,465 & 522 & 3\,457 & 7.5\,s & 1.4\,s \\
\quad full v12b 50k   & 2 & 49\,680 & 1\,602 & 375 & 2\,909 & 6.7\,s & 1.1\,s \\
\midrule
\multicolumn{8}{l}{\textit{n = 1.0\,\AA}} \\
\quad 20k baseline    & 11 & 18\,737 & 3\,595 & 2\,230 & 9\,345 & 14.3\,s & 7.3\,s \\
\quad 50k baseline    & 11 & 48\,278 & 5\,009 & 2\,309 & 12\,557 & 19.2\,s & 7.7\,s \\
\quad full v12b 50k   & 12 & 50\,000 & 4\,958 & 2\,086 & 12\,776 & 17.7\,s & 7.2\,s \\
\midrule
\multicolumn{8}{l}{\textit{n = 1.5\,\AA}} \\
\quad 20k baseline    & 1 & 20\,000 & 7\,281 & 4\,062 & 20\,000 & 21.3\,s & 12.1\,s \\
\quad 50k baseline    & 1 & 50\,000 & 10\,398 & 4\,062 & 35\,690 & 30.3\,s & 12.2\,s \\
\quad full v12b 50k   & 1 & 50\,000 & 6\,067 & 3\,169 & 15\,734 & 21.4\,s & 10.9\,s \\
\midrule
\multicolumn{8}{l}{\textit{n = 2.0\,\AA}} \\
\quad 20k baseline    & 1 & 20\,000 & 3\,494 & 813 & 13\,888 & 10.0\,s & 2.3\,s \\
\quad 50k baseline    & 1 & 50\,000 & 4\,968 & 813 & 13\,888 & 14.0\,s & 2.2\,s \\
\quad full v12b 50k   & 1 & 50\,000 & 3\,219 & 817 & 10\,545 & 11.2\,s & 2.9\,s \\
\bottomrule
\end{tabular}
\end{table}

\begin{table}[h]
\centering\small
\caption{Median steps and wall time: LJ vs.\ DFTB0 (full v12b 50k).}
\begin{tabular}{ccccccc}
\toprule
 & \multicolumn{2}{c}{Median steps} & \multicolumn{2}{c}{Mean wall} & \multicolumn{2}{c}{Speedup} \\
\cmidrule(lr){2-3}\cmidrule(lr){4-5}\cmidrule(lr){6-7}
Noise & LJ & DFTB0 & LJ & DFTB0 & Steps & Wall \\
\midrule
0.5\,\AA & 375  & 2\,613  & 6.7\,s  & 74\,s   & $7\times$ & $11\times$ \\
1.0\,\AA & 2\,086 & 5\,435 & 17.7\,s & 165\,s  & $2.6\times$ & $9\times$ \\
1.5\,\AA & 3\,169 & 7\,755 & 21.4\,s & 238\,s  & $2.4\times$ & $11\times$ \\
2.0\,\AA & 817  & ${\sim}$11k & 11.2\,s & 368\,s  & $14\times$ & $33\times$ \\
\bottomrule
\end{tabular}
\end{table}

% =============================================================================
\section{Gain Decomposition}
\label{app:gains}
% =============================================================================

\begin{table}[h]
\centering\small
\caption{Gain decomposition (strict): LJ vs.\ DFTB0 v13.}
\begin{tabular}{ccccccc}
\toprule
 & \multicolumn{3}{c}{LJ v2} & \multicolumn{3}{c}{DFTB0 v13} \\
\cmidrule(lr){2-4}\cmidrule(lr){5-7}
Noise & $\Delta_\text{50k}$ & $\Delta_\text{kicks}$ & Total
      & $\Delta_\text{50k}$ & $\Delta_\text{kicks}$ & Total \\
\midrule
0.5 & $+$1 ($+$0.3) & $+$7 ($+$2.4) & $+$8
    & $+$3 ($+$1.0) & $+$6 ($+$2.1) & $+$9 \\
1.0 & $+$12 ($+$4.2) & $+$5 ($+$1.7) & $+$17
    & $+$12 ($+$4.2) & $+$14 ($+$4.9) & $+$26 \\
1.5 & $+$22 ($+$7.7) & $+$4 ($+$1.4) & $+$26
    & $+$34 ($+$11.8) & $+$19 ($+$6.6) & $+$53 \\
2.0 & $+$4 ($+$1.4) & $+$2 ($+$0.7) & $+$6
    & $+$50 ($+$17.4) & $+$35 ($+$12.2) & $+$85 \\
\bottomrule
\end{tabular}
\end{table}

Kicks contribute less on LJ at every noise level, because the dominant
failure (ghost modes) cannot be resolved by perturbation.

% =============================================================================
\section{Failure Eigenvalue Spectra}
\label{app:failures}
% =============================================================================

\subsection{Error counts}

Every combo produces 2--14 errors from a pre-existing \texttt{out\_new}
\texttt{UnboundLocalError} bug (code bug, not optimizer failure).

\begin{center}\small
\begin{tabular}{lcccc}
\toprule
Config & $n{=}0.5$ & $n{=}1.0$ & $n{=}1.5$ & $n{=}2.0$ \\
\midrule
All configs (identical) & 14 & 8 & 5 & 2 \\
\bottomrule
\end{tabular}
\end{center}

\subsection{$n{=}0.5$: mild ghost modes}

Full v12b 50k: only 2 failures.  Sample 35:
$\lambda_\text{min} = -6.7 \times 10^{-4}$.
Sample 180: $\lambda_\text{min} = -1.35 \times 10^{-3}$.
Both pass at threshold 0.002.

\subsection{$n{=}1.0$: single-mode failures}

6 failures, all with exactly 1 negative eigenvalue in $(-0.002, 0)$.
All ran 50k steps.  Not real saddle points.

\subsection{$n{=}1.5$: multi-mode cascade}

20k baseline: 53 failures with 1--10 neg modes.  50k: 15.
Full v12b: only 3 (all pass at threshold 0.001).

\subsection{$n{=}2.0$: ghost mode dominance}

Full v12b 50k autopsy: ghost modes 191/287 (67\%), drifting 38 (13\%),
oscillating 33 (12\%), strict converged 19 (7\%).
Only 1 non-converged failure (sample 178, 2 neg modes, 50k exhausted).

% =============================================================================
\section{Ablation Results at $n{=}2.0$}
\label{app:ablations}
% =============================================================================

\begin{table}[h]
\centering\small
\caption{Ablations at $n{=}2.0$\,\AA.  Converged includes relaxed-accept.}
\begin{tabular}{lccccc}
\toprule
Configuration & Strict & Converged & Non-conv & Med.\ steps & Wall \\
\midrule
20k baseline (no kicks) & 13 (4.5\%) & 263 & 22 & 813 & 10\,s \\
50k baseline (no kicks) & 17 (5.9\%) & 275 & 10 & 813 & 14\,s \\
\midrule
Full v12b, 20k          & 16 (5.6\%) & 276 & 9 & 817 & 10\,s \\
Full v12b, 50k          & 19 (6.6\%) & 284 & 1 & 817 & 11\,s \\
Adapt+probe (no late)   & 20 (7.0\%) & 284 & 1 & 817 & 11\,s \\
\textbf{Probe+late (no adapt)} & \textbf{21} (7.3\%) & \textbf{285} & \textbf{0} & 866 & 11\,s \\
\bottomrule
\end{tabular}
\end{table}

Key findings: (1)~step count matters more than kicks, (2)~late escape
unnecessary at 50k, (3)~adaptive scaling unnecessary on smooth LJ,
(4)~all 50k configs give 284--285/287 under relaxed---differences are
ghost-mode noise, not meaningful.

% =============================================================================
\section{Hardest Samples}
\label{app:hardest}
% =============================================================================

\begin{center}\small
\begin{tabular}{ccc|ccc}
\toprule
\multicolumn{3}{c}{\textbf{\color{ljblue}LJ hardest}} &
\multicolumn{3}{c}{\textbf{\color{dftborange}DFTB0 hardest}} \\
Sample & Conv & Rate & Sample & Conv & Rate \\
\midrule
106 & 4/15 & 27\% & 136 & 6/15 & 40\% \\
35  & 6/15 & 40\% & 150 & 6/15 & 40\% \\
178 & 8/15 & 53\% & 209 & 7/15 & 47\% \\
71  & 9/15 & 60\% & 212 & 8/15 & 53\% \\
67  & 10/15 & 67\% & 51  & 9/15 & 60\% \\
\bottomrule
\end{tabular}
\end{center}

\textbf{Zero overlap.}  Difficulty is PES-specific, not geometry-specific.

% =============================================================================
\section{Theoretical Discussion}
\label{app:theory}
% =============================================================================

\subsection{Validation of the $\varepsilon^{1/2}$ theorem}

The $\varepsilon^{1/2}$ theorem (Boumal et al.) predicts
$\lambda_\text{min} \geq -C\sqrt{\varepsilon}$ for RFO-type methods.
With $\varepsilon = 10^{-4}$: $\lambda_\text{min} \geq -0.01$.

LJ provides striking validation: at $n{=}2.0$, relaxing to $-0.01$
rescues 285/287 (99.3\%).  At $n{=}1.0$, all 6 failures have eigenvalues
in $(-0.002, 0)$.  The cascade knee at $10^{-4}$--$10^{-3}$ confirms
the theorem's bound is conservative.

On DFTB0 the theorem is invisible (strict $\approx$ relaxed).  LJ
exposes it by providing genuinely flat directions.

\subsection{Pairwise rank deficiency}

For $V = \sum_{i<j} V(r_{ij})$, pure angular deformations (bond rotations
about the axis connecting terminal atoms) change no pairwise distance to
first order.  These zero-eigenvalue directions are \emph{intrinsic to the
potential}.  At equilibrium, coupling through the pairwise network may
give quasi-angular modes positive curvature, but noise displaces the
geometry into flat regions, increasing the number of near-zero eigenvalues.

\subsection{Cascade gap scaling}

\begin{center}\small
\begin{tabular}{cccc}
\toprule
Noise $\sigma$ & LJ gap (pts) & DFTB0 gap & Ratio \\
\midrule
0.5\,\AA & 2.8 & ${\sim}$0 & --- \\
1.0\,\AA & 32.7 & ${\sim}$0 & ${\sim}\infty$ \\
1.5\,\AA & 64.9 & ${\sim}$0 & ${\sim}\infty$ \\
2.0\,\AA & 92.7 & ${\sim}$0 & ${\sim}\infty$ \\
\bottomrule
\end{tabular}
\end{center}

Gap scales roughly as $\sigma^2$, consistent with the number of flat
directions growing with displacement magnitude.

% =============================================================================
\section{Known Issues}
\label{app:issues}
% =============================================================================

\begin{enumerate}[nosep]
  \item \textbf{\texttt{out\_new} bug:}  5--14 samples/combo crash with
    \texttt{UnboundLocalError} in the trust-radius retry loop.  Code bug,
    not optimizer limitation.
  \item \textbf{Seed consistency:}  Different seeds produce identical
    results at $n{=}2.0$ full v12b 50k (strict=19, relaxed=285, same
    single failure sample 178).
  \item \textbf{Hessian precision:}  LJ Hessians computed in float64 to
    avoid \texttt{linalg.eigh} failures from wide eigenvalue spectrum.
\end{enumerate}

% =============================================================================
% Data reference
% =============================================================================
% Grid output:  /scratch/memoozd/ts-tools-scratch/runs/min_nr_lj_v2_grid/
% Analysis:     /scratch/memoozd/ts-tools-scratch/runs/min_nr_lj_v2_grid/analysis/
% SLURM script: src/noisy/v2_tests/scripts/slurm_templates/minimization_nr_lj_v2_run.slurm
% Seed: 42, 48 workers x 4 threads, 192 CPUs, completed 2026-03-23 17:41 EDT

\end{document}
